\documentclass[11pt]{article}

\usepackage[margin=2.5cm]{geometry}
\usepackage{amsmath, amssymb}
\usepackage{enumitem}
\setlist{nosep, leftmargin=*, topsep=2pt, itemsep=2pt}
\usepackage{hyperref}
\hypersetup{colorlinks=true, linkcolor=blue!50!black, urlcolor=blue!50!black,
            pdftitle={Outline -- Rewiring the market time clock in Spanish electricity markets},
            pdfauthor={Pablo Paramio Perez}}
\usepackage{microtype}
\linespread{1.10}

\title{Outline\\
\large Rewiring the market time clock in Spanish electricity markets}
\author{Pablo Paramio P\'erez}
\date{\today}

\begin{document}
\maketitle

\section*{RAP}

\noindent\textbf{Positioning.} How shorter trading intervals reshape
wholesale electricity bidding and prices has been studied for Germany
(M\"arkle-Hu\ss, Feuerriegel and Neumann, 2017) and Australia
(Csereklyei and Khezr, 2024). The recent two-step Spanish move to
15-minute trading --- intraday auctions in March 2025, day-ahead
market in October 2025 --- has not yet been evaluated, and it lands
in a setting reshaped by the April 2025 blackout and the continuous
reinforced-operation regime that followed, requiring a careful
empirical strategy to disentangle the reform's effect from the
regulatory response to the blackout.

\noindent\textbf{Research Question.} Does shorter trading-period
granularity reshape the strategic margin in sequential electricity
markets? Who uses the within-hour freedom, when during the day is it
used, through which equilibrium channel does the reform operate, and
does it amplify or compress market power? Do prices fall or rise,
and how does the price gap between sequential markets evolve?

\noindent\textbf{Answer.} Granularity acts as a \emph{strategic-latitude
reveal}: it does not change marginal costs, capacity or who
participates; it expands the room for within-hour bid discrimination,
and the equilibrium response traces exactly that.

We build a sequential-market Cournot model with a strategic bidder, a
forecast-driven renewable aggregator, and a capacity-bound arbitrageur
linking the sequential markets, and show that the reform expands
within-hour bid-ladder discrimination, pushes prices down on the side
of the market that becomes finer (with the other market co-moving
down in anticipation), and shifts forecast-driven volumes into the
finer market while the average price gap between sequential markets
stays at zero. The arbitrageur's bounded capacity prevents the
within-hour price-gap dispersion from collapsing, so within-hour
heterogeneity persists across both reform steps, and the welfare
reading is consumer-favourable on the auction side.

We confirm these mechanisms in rich Spanish market and bid-level data.
Conventional generators widen their within-hour bidding schedules
where within-hour residual demand actually varies; prices on the side
of the market that becomes finer fall sharply with the other market
following down in anticipation; more conventional generation clears
after the day-ahead reform; the within-hour price-gap dispersion
jumps under the first reform and does not collapse at the second;
and the intraday bid stack itself reveals active position rebalancing
by producers across every technology. The reform therefore compresses
rather than amplifies market power, and the post-blackout cost rise
on the same calendar window is a separate regulatory response, not a
granularity effect.


\section*{Questions Provoked by the RAP}

\begin{enumerate}
  \item Which equilibrium channel of the sequential-market model
        carries the reform's effect, and which firm types respond
        through which channel?
  \item How is the day-ahead cutover identified given the concurrent
        post-blackout reforzada regime?
  \item Does the reform amplify market power, and does it harm
        consumers, once both the wholesale auction and the system-cost
        cascade are accounted for?
  \item Could the March 2025 intraday reform have contributed to the
        macro-conditions that made the April 2025 blackout possible?
\end{enumerate}

\section{Introduction}

The introduction motivates the research question against the German and Australian precedents on shorter trading intervals, states the RAP, and previews the headline take-away --- granularity is a strategic-latitude reveal, not a market-power amplifier.\hfill\emph{Target: $\sim 1.5$ pp.}

\section{Institutional setting: sequential markets, three reforms, and a blackout}

This section maps the Spanish wholesale market layer and describes the three 15-minute reforms and the April 2025 blackout that switches the system into the continuous reinforced-operation regime. Key fact: the reinforced-operation regime overlaps only with the day-ahead reform, which forces the empirical strategy to hold it constant by design.\hfill\emph{Target: $\sim 2$ pp.}
\\\emph{Skeleton:}
\begin{itemize}
  \item \textbf{2.1} Spanish wholesale architecture: sequential auctions and balancing.
  \item \textbf{2.2} The three 15-minute reforms (settlement ISP15, intraday ID15, day-ahead DA15).
  \item \textbf{2.3} The April 2025 blackout and the reinforced-operation regime.
\end{itemize}

\section{Theory: a sequential-market Cournot model with granularity selectors}

This section builds the two-bidder sequential-market Cournot model with three agents (strategic bidder, operational aggregator, capacity-bound arbitrageur) and a granularity selector at each stage, traces the equilibrium through the three reform states, derives the seven sign predictions, and contrasts the empirically relevant limited-arbitrage regime with the full and strategic benchmarks. Key result: all welfare channels are sign-positive under uniform within-hour shocks, and the empirical asymmetric persistence of within-hour price-gap dispersion is the equilibrium signature of limited per-product arbitrage, not a welfare cost. Closed-form derivations in Appendix~D.\hfill\emph{Target: $\sim 6$ pp.}
\\\emph{Skeleton:}
\begin{itemize}
  \item \textbf{3.1} Setup: three sequential auctions, three agents.
  \item \textbf{3.2} Arbitrage regimes: four benchmarks, one empirical.
  \item \textbf{3.3} Pre-reform $(1,1)$: scarcity wedge, no within-hour dispersion.
  \item \textbf{3.4} Asymmetric state $(1,Q)$: four spot-side mechanisms.
  \item \textbf{3.5} Symmetric state $(Q,Q)$: limited per-product arbitrage.
  \item \textbf{3.6} Comparative statics across the reform path.
  \item \textbf{3.7} Welfare across reforms and arbitrage regimes.
\end{itemize}

\section{Data: OMIE bid-level offers, weather, redispatch}

This section describes the OMIE bid-level offers (every unit, date and period in the day-ahead and intraday auctions, 2022 to 2026), the 15-minute ENTSO-E weather and load series, and the ESIOS redispatch volumes and prices; defines the five outcome families used downstream; explains the window-and-market-specific bandwidth that isolates the competing tier; and closes with the critical/flat partition of clock-hours --- the within-day classifier that organises the empirical strategy.\hfill\emph{Target: $\sim 1.5$ pp.}
\\\emph{Skeleton:}
\begin{itemize}
  \item \textbf{4.1} Sources: OMIE bid-level offers, ENTSO-E weather and load, ESIOS redispatch.
  \item \textbf{4.2} Outcomes by family: bid-shape, in-band quantity, system aggregates, cross-market wedge, within-hour dispersion.
  \item \textbf{4.3} Window-and-market-specific bandwidth and the in-band region.
  \item \textbf{4.4} The critical/flat partition: where granular pricing has economic content.
\end{itemize}

\section{Empirical strategy: three specifications for three economic objects}

This section lays out the three identification specifications --- per-curve difference-in-differences on bid shape, per-hour-class Bayesian counterfactual on in-band offered energy, daily Bayesian counterfactual on prices and per-technology cleared MW --- and explains how the day-ahead reform is identified holding the reinforced-operation regime constant by starting the pre-window on the day of the blackout.\hfill\emph{Target: $\sim 4$--$5$ pp.}
\\\emph{Skeleton:}
\begin{itemize}
  \item \textbf{5.1} Spec A: Bayesian Structural Time Series for prices, cleared energy, and wedge SD.
  \item \textbf{5.2} Spec B: per-hour-class BSTS on bid-level outcomes.
  \item \textbf{5.3} Spec C: per-curve DiD for bid shape.
\end{itemize}

\section{Results}

This section presents the empirical findings organised one per model prediction. Take-away: dispatchable firms widen their within-quarter bid ladders in critical hours, intraday clearing prices fall sharply with the day-ahead price following down in anticipation, day-ahead combined-cycle gas cleared volume rises at the day-ahead reform, within-hour price-gap dispersion jumps under the intraday reform and does not collapse at the day-ahead reform, and the intraday bid stack reveals an active arbitrage layer across every technology that is actually clearing volume.\hfill\emph{Target: $\sim 10$ pp.}
\\\emph{Skeleton:}
\begin{itemize}
  \item \textbf{6.1} Bid-shape widening in critical hours --- both reforms (Spec C).
  \item \textbf{6.2} In-band offered energy migrates into the granular market --- both reforms (Spec B).
  \item \textbf{6.3} ID15 clearing-price drop in both markets via cross-market anticipation (Spec A).
  \item \textbf{6.4} DA15 day-ahead combined-cycle gas cleared scale-up in critical hours (Spec A).
  \item \textbf{6.5} Within-hour wedge SD jumps under ID15 and persists at DA15 (Spec A).
  \item \textbf{6.6} Production-unit arbitrageurs are empirically active in intraday auctions.
  \item \textbf{6.7} Cross-market substitution and operational-vs-strategic incentives explain the wind null.
\end{itemize}

\section{Conclusion: granularity is a strategic-latitude reveal, not a market-power amplifier}

The conclusion synthesises the auction-side, system-cost and bid-side evidence against the model's welfare reading; reads the reform as a strategic-latitude reveal rather than a market-power amplifier; separates the post-blackout reinforced-operation cost rise from the granularity reforms; and flags the research gap on whether the March 2025 intraday reform contributed to the macro-conditions of the April 2025 blackout, beyond the scope of this thesis.\hfill\emph{Target: $\sim 3$ pp.}

\section*{Appendix A: Robustness}

This appendix collects the standard robustness checks for the main-body findings --- within-pre placebos, hour-class falsifications, cross-border-flow control, per-firm and per-session decompositions, bandwidth sensitivity, and offer-type composition. Take-away: all headline results survive every check, and the simple-vs-block CCGT composition shift dampens rather than drives the bid-shape widening.\hfill\emph{Target: $\sim 2$ pp.}
\\\emph{Skeleton:}
\begin{itemize}
  \item \textbf{A.1} Within-pre midpoint placebos for bid-shape DiDs.
  \item \textbf{A.2} Midday-vs-flat hour-class falsification.
  \item \textbf{A.3} Cross-border-flow control on intraday $\sigma_p$.
  \item \textbf{A.4} Per-CCGT-firm decomposition of the DA15 $\sigma_p$ widening.
  \item \textbf{A.5} Bandwidth sensitivity (the $p_{90}$ choice and wider alternatives).
  \item \textbf{A.6} Per-IDA-session BSTS for the intraday price drop.
  \item \textbf{A.7} Offer-type composition check (simple / MIC / block).
\end{itemize}

\section*{Appendix B: Continuous-market response at the three reforms}

This appendix runs the Spec A BSTS architecture on the continuous-intraday market separately for each of the three reform cutovers. Take-away: a methodological null at the settlement-only reform; a structural response at the intraday cutover (trade count quadruples, energy turnover drops as smaller contracts cannibalise, prices follow the auction-side drop); a smaller indirect effect at the day-ahead cutover.\hfill\emph{Target: $\sim 2$ pp.}
\\\emph{Skeleton:}
\begin{itemize}
  \item \textbf{B.1} ISP15 (settlement-only): a methodological null.
  \item \textbf{B.2} ID15 ($60\to15$ min): the structural continuous-market reform.
  \item \textbf{B.3} DA15: a smaller, indirect continuous-market effect.
\end{itemize}

\section*{Appendix C: Market-structure context}

This appendix presents the market-structure facts that frame the bidding evidence: the two-tier combined-cycle gas day-ahead curve, the ancillary-market concentration migration, the granularity-vs-reinforced-operation cost decomposition, and the direct evidence on unit-level arbitrage. Take-away: market power has migrated downstream of the wholesale clear (HHI moderate at wholesale, high in ancillary services); the granularity reforms shrink the auction's adjustment-cost channels while the reinforced-operation regime drives the post-blackout cost rise on combined-cycle gas only.\hfill\emph{Target: $\sim 5$ pp.}
\\\emph{Skeleton:}
\begin{itemize}
  \item \textbf{C.1} Two-tier combined-cycle gas day-ahead curve and the day-ahead/intraday contrast.
  \item \textbf{C.2} Market concentration migrated from wholesale to ancillary services.
  \item \textbf{C.3} Granularity shrinks adjustment-cost channels; reinforced operation inflates Fase~I.
  \item \textbf{C.4} The reinforced-operation cost increase concentrates on combined-cycle gas.
  \item \textbf{C.5} Production-unit arbitrageurs are empirically active in the intraday auctions.
\end{itemize}

\section*{Appendix D: Theoretical derivations}

This appendix provides the full closed-form derivations underlying the model in Section~3, organised by state $(1,1) \to (1,Q) \to (Q,Q)$ and by agent (operational aggregator scheduling problem, strategic bidder two-tranche ladder, capacity-bound arbitrageur). Take-away: all comparative statics admit closed forms under uniform within-hour shocks, and the per-step welfare delta is unconditionally positive on the imbalance and Cournot channels.\hfill\emph{Target: $\sim 11$ pp.}
\\\emph{Skeleton:}
\begin{itemize}
  \item \textbf{D.1} Setup details and information arrival.
  \item \textbf{D.2} Operational aggregator: full scheduling problem.
  \item \textbf{D.3} Strategic bidder: optimal two-tranche ladder.
  \item \textbf{D.4} $(1,1)$ equilibrium across the four arbitrage regimes.
  \item \textbf{D.5} $(1,Q)$ equilibrium: block parity, strategic arbitrage, wedge SD.
  \item \textbf{D.6} $(Q,Q)$ equilibrium: per-product limited arbitrage and Cournot scale-up.
  \item \textbf{D.7} Balancing discrepancy across the three states.
  \item \textbf{D.8} Note on the alternative sequence $(Q,1)$.
  \item \textbf{D.9} Welfare derivations and net per-step welfare deltas.
\end{itemize}

\end{document}
